\scanpage{22}
\uline{\S1} Pseudo-anneau $=$ groupe abélien à loi de composition interne associative bilinéaire.

Nous sous-entendrons par la suite la commutativité.

Pseudo-anneau \inserted{$J$} \uline{à puissances} \uline{divisées}. C'est un ps-anneau muni d'applications \deleted{satisfaisant}
\[
\pi^n:x\rightsquigarrow x^{(n)}=\pi^n(x)\qquad n\in\mathbb N^+
\]
satisfaisant les conditions
\[
\begin{aligned}
1.1.\quad&\pi^1=\mathrm{id}\qquad\text{i.e.}\qquad x^{(1)}=x\\
1.2.\quad&\pi^n(x+y)=\deleted{\unsure{\sum}}\pi^n(x)+\sum_{\substack{p,q\ge1\\p+q=n}}\pi^p(x)\pi^q(y)+\pi^n(y)\\
1.3.\quad&\pi^n(xy)=x^n\pi^n(y)\\
1.4.\quad&\pi^m(\pi^n(x))=\frac{(mn)!}{(n!)^m m!}\pi^{mn}(x)\\
1.5.\quad&\pi^m(x)\pi^n(x)=\frac{(m+n)!}{m!n!}\pi^{m+n}(x)
\end{aligned}
\]
Supposons que $J$ est un idéal d'un anneau $A$, avec la structure de pseudo-anneau induite. Alors la relation 2. peut s'écrire
\[
\pi^n(x+y)=\sum_{p+q=n}\pi^p(x)\pi^q(y)\qquad x,y\in J
\]
où on pose $\pi^0x=1$, $\pi^p(x)=0$ si $p<0$. De la formule 5. on tire
\[
\text{(1.5.bis)}\qquad\pi^{p_1}(x)\pi^{p_2}(x)\cdots\pi^{p_r}(x)=\frac{(p_1+\cdots+p_r)!}{p_1!\cdots p_r!}\pi^{p_1+\cdots+p_r}(x)
\]
en particulier
\[
\text{1.\unsure{5|7}.ter}\qquad n!\pi^n(x)^{\deleted{\unsure{n}}}=\deleted{\unsure{n!}}x^n
\]
\unsure{Que|Si} $\Lambda$ est un anneau (comm), une $\Lambda$-algèbre
\mspage{21}
